% =================================================
\chapter{Pore-pressure experiment: interpretation and inverse analyses}
\label{ch:inversion}

The mine-by pore-pressure measurements performed during the Canadian \acrshort{TSX} experiment motivated a similar study at \URFBII. 
The aim was to estimate in-situ, rock-mass-scale hydro-mechanical parameters from continuous pore-pressure monitoring and modelling, as a complement to conventional laboratory measurements on small core samples.

Initial stresses and hydraulic conditions at \URFB were favourable. Compared with \acrshort{TSX}, a faster pressure relaxation was expected due to the higher hydraulic conductivity; this was compensated for by continuous automated monitoring. At the same time, the geological setting is more complex and the placement of pressure-monitoring boreholes was constrained by technical limitations. These drawbacks were expected to be offset by a larger number of measurements and by the use of full 3D hydro-mechanical models.

A basic description of the experiment and a first interpretation of the pressure series are given in Section~\ref{sec:exp_interpretation}. The measurements clearly demonstrate high sensitivity to excavation and other perturbations; however, the explanation provided by the homogeneous 3D model is insufficient for a meaningful full inversion of hydro-mechanical parameters. We therefore employ two complementary approaches.


First, in Section~\ref{sec:exp_interpretation}, we use Bayesian inversion of the pressure response to \gls{WPT}s in combination with a simple 1D model of an individual measurement cell. The main goals are (i) to obtain a more precise estimate of the initial pressure field prior to excavation and (ii) to provide a basic comparison of hydraulic conductivity and Young’s modulus before and after excavation. Second, in Section~\ref{3D_excavation_model}, we use geological evidence from fracture outcrops on excavation walls and in borehole logs to construct a limited set of candidate discrete fracture surfaces, which are then incorporated into the 3D hydraulic model to improve agreement with the initial pressure field.

Finally, the computational cost of the full 3D model prevents routine sampling-based inversion. To address this limitation, we investigate a delayed-acceptance framework accelerated by neural-network surrogate models in Section~\ref{sec:TSX_heterogeneity}. Since a suitably calibrated 3D Bukov model is not yet available, the surrogate workflow is developed and tested on a benchmark hydro-mechanical model motivated by the \acrshort{TSX} experiment.





% \sts{[Standa: Chybi nejaky uvod k teto kapitole, ktery by dal do souvislosti Sekce 6.1, 6.2 a 6.3. Hlavne Sekce 6.2 momentalne vubec nenavazuje na 6.1 a 6.3. Doporucil bych ji prohodit se Sekci 6.2.]}

\section{Basic interpretation of the pore pressure measurements}
\label{sec:exp_interpretation}
% author Jan Březina
\input 06_pore_pressure_experiment
%
%
\section{Preparation of 3D heterogeneous model}
\label{3D_excavation_model}
% author: Jan Stebel
\input 06_fractures

\section{Accelerated Bayesian techniques}
\label{sec:TSX_heterogeneity}
% author: Michal Béreš
\input 06_TSX_heterogeneity

\section{Chapter summary}
Pore pressure monitoring during the excavation revealed clear responses across all sensors, though some short-term anomalies remain unexplained. Measured pressure changes were generally smaller than predicted by the homogeneous model, highlighting the importance of fractures and heterogeneous hydraulic pathways. Differences between the northern and southern chambers suggest variations in conductivity and rock stiffness.

Initial \Gls{WPT}s were influenced by preceding drilling and insufficient equilibration, emphasizing the need for careful pre-conditioning, stepwise \Gls{WPT} with flux measurements, long-term pressure relaxation, and controlled return to equilibrium before excavation.

A heterogeneous hydro-mechanical model incorporating explicit fracture planes substantially improved agreement with observations, reducing the pressure misfit to about $9~\%$ of the homogeneous model value, though some hydraulically isolated zones remain unrepresented.

Bayesian inversion using delayed acceptance and adaptive surrogate techniques efficiently identified heterogeneity in hydraulic and elastic properties, demonstrating the critical role of fractures and structural variability in controlling pore pressure responses.
